\chapter*{Introduction}
\addcontentsline{toc}{chapter}{Introduction}

Dans le cadre de ce projet intégratif, nous étions chargés de concevoir un robot capable de répondre à une problématique imposée. Ce projet avait pour objectif de mettre en \oe uvre les différentes compétences acquises lors du cours d'Algorithmique ainsi que celui d'Électronique pour les systèmes embarqués. La gestion de projet était également importante pour assurer le suivi et la coordination des différentes étapes du projet.\\


\noindent Ce projet se divise en deux parties complémentaires. La première partie est une partie algorithmique, qui consiste en l'analyse de la description de la ville. Afin de déterminer le parcours le plus optimal, et préparer une liste d'instructions pour le robot. La deuxième partie est une partie robotique, qui consiste à concevoir le robot, et faire un programme capable de le piloter, et d'exécuter les différentes instructions générées lors de la première partie.\\

\noindent Le cahier des charges pour le projet sera le suivant :
\begin{itemize}
    \item Le robot doit être capable de suivre des lignes noires et de corriger sa trajectoire en cas de déviations.
    \item Le robot doit être capable de franchir un certain nombre de points de passage imposés.
    \item Le robot doit être capable de revenir à son point de départ.
    \item L'algorithme doit être capable de trouver le parcours le plus court possible.
\end{itemize}

\vspace{1cm}

\noindent Le rapport sera donc structuré de la manière suivante. Dans une première partie, nous présenterons, de manière plus détaillée les contraintes techniques du sujet, l'organisation du projet et les différentes étapes de conception. Dans une seconde partie, nous présenterons la conception électronique du robot, incluant le schéma électrique, les choix techniques et le bilan de puissance du robot. Ensuite, nous aborderons dans la troisième partie, les choix de conception et d'implémentation de notre robot. Celui-ci suivra le cycle en V, et abordera les algorithmes clés. Dans un dernier temps, nous ferons un bilan de la démonstration, avec un commentaire sur les performances du robot, et d'éventuels problèmes rencontrés. Enfin, nous ferons une synthèse du projet, en abordant les différents problèmes rencontrés et les solutions proposées à une échelle plus globale.\\
